\documentclass[11pt,letterpaper]{article}
\usepackage[T1]{fontenc}
\usepackage[scaled]{helvet}
\renewcommand{\familydefault}{\sfdefault}
\usepackage{amsmath,amssymb}
\usepackage[margin=0.8in,bottom=0.9in]{geometry}
\usepackage[most]{tcolorbox}
\usepackage{tikz}
\usetikzlibrary{calc}
\usepackage{enumitem}
\usepackage{array}
\usepackage{needspace}
\usepackage{fancyhdr}
\usepackage{pifont}

% ---- Palette "F. Bright & Playful" ----
\definecolor{fmtitle}{HTML}{3D2C8D}
\definecolor{fmkey}{HTML}{0096C7}
\definecolor{fmwatch}{HTML}{2BB673}
\definecolor{fmturn}{HTML}{F4A100}
\definecolor{fmmistake}{HTML}{EF476F}
\definecolor{fmnotice}{HTML}{9B5DE5}
\definecolor{fmgray}{HTML}{6B6B6B}

\pagestyle{fancy}
\fancyhf{}
\renewcommand{\headrulewidth}{0pt}
\fancyfoot[C]{\footnotesize\color{fmgray}\textcopyright\ 2026 FangMath. All rights reserved.}
\fancyfoot[R]{\footnotesize\color{fmgray}\thepage}
\setlength{\parindent}{0pt}
\setlength{\parskip}{3pt}

\newcommand{\badge}[1]{\colorbox{fmtitle}{\color{white}\bfseries\small\ #1\ }}
\newcommand{\leveltag}[2]{\colorbox{#1}{\color{white}\bfseries\scriptsize\ #2\ }}
\newcommand{\LE}{\leveltag{fmwatch}{EASY}}
\newcommand{\LM}{\leveltag{fmturn}{MEDIUM}}
\newcommand{\LH}{\leveltag{fmmistake}{HARD}}

\newcommand{\sectionbar}[2]{\needspace{6cm}\par\vspace{4mm}
  \begin{tcolorbox}[enhanced,colback=fmtitle,colframe=fmtitle,arc=2mm,
    left=4mm,top=1.5mm,bottom=1.5mm]
  \color{white}\bfseries\large #1\hfill{\normalsize\mdseries #2}
  \end{tcolorbox}}

\newtcolorbox{keyidea}{enhanced,breakable,colback=white,colframe=fmkey,boxrule=1pt,arc=2mm,
  title={Key Idea},coltitle=fmkey,fonttitle=\bfseries,
  attach boxed title to top left={xshift=4mm,yshift=-2mm},
  boxed title style={colback=white,colframe=fmkey,boxrule=0.9pt}}
\newtcolorbox{watchme}{enhanced,colback=white,colframe=fmwatch,boxrule=1pt,arc=2mm,
  title={Watch Me --- Worked Example},coltitle=fmwatch,fonttitle=\bfseries,
  attach boxed title to top left={xshift=4mm,yshift=-2mm},
  boxed title style={colback=white,colframe=fmwatch,boxrule=0.9pt}}
\newtcolorbox{mistakebox}{enhanced,breakable,colback=white,colframe=fmmistake,boxrule=1pt,arc=2mm,
  title={Common Mistake},coltitle=fmmistake,fonttitle=\bfseries,
  attach boxed title to top left={xshift=4mm,yshift=-2mm},
  boxed title style={colback=white,colframe=fmmistake,boxrule=0.9pt}}
\newtcolorbox{yourturn}{enhanced,breakable,colback=white,colframe=fmturn,boxrule=1pt,arc=2mm,
  title={Your Turn},coltitle=fmturn,fonttitle=\bfseries,
  attach boxed title to top left={xshift=4mm,yshift=-2mm},
  boxed title style={colback=white,colframe=fmturn,boxrule=0.9pt}}
\newtcolorbox{answerbox}[1]{enhanced,breakable,colback=white,colframe=fmwatch,boxrule=1pt,arc=2mm,
  title={#1},coltitle=fmwatch,fonttitle=\bfseries,
  attach boxed title to top left={xshift=4mm,yshift=-2mm},
  boxed title style={colback=white,colframe=fmwatch,boxrule=0.9pt}}

\newcommand{\work}[1]{\par\vspace{#1}}
\setlist[enumerate]{itemsep=2pt,topsep=3pt}

% ---------------- extra boxes ----------------
\newtcolorbox{notebox}{enhanced,breakable,colback=white,colframe=fmnotice,boxrule=1pt,arc=2mm,
  title={Notice},coltitle=fmnotice,fonttitle=\bfseries,
  attach boxed title to top left={xshift=4mm,yshift=-2mm},
  boxed title style={colback=white,colframe=fmnotice,boxrule=0.9pt}}
\newtcolorbox{exbox}{enhanced,colback=white,colframe=fmwatch,boxrule=0.8pt,arc=2mm,
  title={Example},coltitle=fmwatch,fonttitle=\bfseries\small,
  attach boxed title to top left={xshift=4mm,yshift=-2mm},
  boxed title style={colback=white,colframe=fmwatch,boxrule=0.8pt},fontupper=\small}
\newcommand{\topicbar}[1]{\needspace{4.5cm}\par\vspace{5mm}
  \begin{tcolorbox}[enhanced,colback=fmtitle,colframe=fmtitle,arc=2mm,left=4mm,top=1.5mm,bottom=1.5mm]
  \color{white}\bfseries\large #1
  \end{tcolorbox}}
\usepackage{multicol}

% ---------------- practice layout ----------------
\newcommand{\workline}{\par\vspace{7mm}{\color{gray!55}\rule{\linewidth}{0.4pt}}}
\newcommand{\worklines}[1]{\foreach \i in {1,...,#1}{\workline}}
\newcommand{\prob}[3]{\needspace{#3}\par\vspace{3mm}{\color{fmturn}\bfseries #1.}\ #2\par}
\newcommand{\ansline}{\par\vspace{4mm}\hfill\textbf{Answer:}\ \rule{5cm}{0.4pt}\par}
\newcommand{\keyitem}[3]{\needspace{3.2cm}\par\vspace{2mm}
  \begin{tcolorbox}[enhanced,colback=white,colframe=fmwatch!60,boxrule=0.6pt,arc=1.5mm,left=3mm,top=1.5mm,bottom=1.5mm]
  {\color{fmwatch}\bfseries #1.}\quad\small #2\par\vspace{1mm}
  \normalsize\textbf{Answer: #3}
  \end{tcolorbox}}


\usetikzlibrary{arrows.meta}
\providecommand{\iskey}{0}
% fill-in blank: line for students, pink answer in the key
\newcommand{\fb}[2][2.6cm]{\ifnum\iskey=1 \underline{\makebox[#1][c]{\color{fmmistake}\bfseries #2}}\else\underline{\makebox[#1]{}}\fi}
% answer line for practice problems
\newcommand{\pans}[2][3.2cm]{\ifnum\iskey=1 \underline{\makebox[#1][c]{\color{fmmistake}\bfseries #2}}\else\rule{#1}{0.4pt}\fi}
\newcommand{\vid}[1]{{\color{fmgray}\small\textit{\ding{43} Video: #1}}}
\newcommand{\pt}[2]{\fill (#1) circle (2.2pt) node[above=1pt,font=\small]{$#2$};}

\usetikzlibrary{decorations.pathreplacing}
% three points on a segment with lengths written above the pieces: #1,#2,#3 letters; #4,#5 piece labels; #6 whole label (below)
\newcommand{\threept}[6]{\begin{tikzpicture}[x=1cm]\draw[thick] (0,0)--(3.6,0);
\fill (0,0) circle(2.2pt) node[below=2pt,font=\small]{$#1$} (1.4,0) circle(2.2pt) node[below=2pt,font=\small]{$#2$} (3.6,0) circle(2.2pt) node[below=2pt,font=\small]{$#3$};
\node[above,font=\small] at (0.7,0.02){#4};\node[above,font=\small] at (2.5,0.02){#5};
\if\relax\detokenize{#6}\relax\else\draw[gray,decorate,decoration={brace,amplitude=4pt,mirror}] (0,-0.45)--(3.6,-0.45) node[midway,below=4pt,font=\small]{#6};\fi
\end{tikzpicture}}
\newcommand{\sdp}[3]{\begin{minipage}[t][2.4cm]{0.48\linewidth}{\color{fmturn}\bfseries #1.}\ #2\par\vfill\hfill \pans[3.6cm]{#3}\end{minipage}}
\newcommand{\figp}[4]{\begin{minipage}[t]{0.48\linewidth}{\color{fmturn}\bfseries #1.}\ #2\par\vspace{1mm}\centering #3\par\vspace{2mm}\raggedleft \pans[2.6cm]{#4}\end{minipage}}
\newcommand{\algp}[3]{\begin{minipage}[t][4.6cm]{0.48\linewidth}{\color{fmturn}\bfseries #1.}\ #2\par\vfill\hfill \pans[4.2cm]{#3}\end{minipage}}

\begin{document}
\ifnum\iskey=1 \colorbox{fmwatch}{\color{white}\bfseries\small\ ANSWER KEY\ }\else\badge{PRACTICE}\fi\hfill{\small\color{fmgray}Geometry \textperiodcentered\ Lesson 1.2}\par\vspace{4mm}
{\color{fmtitle}\LARGE\bfseries Segment Addition Postulate --- Practice\ifnum\iskey=1 \ --- Answer Key\fi}\par\vspace{1mm}
{\color{fmgray}A worked example comes before each type. In every problem, the middle point is between the other two unless it says otherwise. Show your thinking.}\par\vspace{3mm}
Name: \rule{6cm}{0.4pt}\hspace{1cm}Date: \rule{3.5cm}{0.4pt}\par\vspace{2mm}

\needspace{6cm}\topicbar{Type 1 \quad Find the Whole: Add}
\begin{watchme}\begin{minipage}[c]{0.4\linewidth}\centering\threept{A}{B}{C}{$5$}{$9$}{$AC = \,?$}\end{minipage}\hfill
\begin{minipage}[c]{0.56\linewidth}$B$ is between $A$ and $C$, so part $+$ part $=$ whole.\par
$AC = AB + BC = 5 + 9$. \quad \textbf{Answer: $\boldsymbol{AC = 14}$}\end{minipage}\end{watchme}
\needspace{3.4cm}\noindent\sdp{1}{$AB = 12$ and $BC = 15$. Find $AC$.}{$AC = 27$}\hfill\sdp{2}{$AB = 3.5$ and $BC = 6.2$. Find $AC$.}{$AC = 9.7$}\par\vspace{2mm}
\needspace{3.4cm}\noindent\sdp{3}{$AB = 2\tfrac12$ and $BC = 4\tfrac14$. Find $AC$.}{$AC = 6\tfrac34$}\hfill\figp{4}{Find $RT$.}{\threept{R}{S}{T}{$8$}{$13$}{$RT = \,?$}}{$RT = 21$}\par

\needspace{6cm}\topicbar{Type 2 \quad Find a Missing Piece: Subtract}
\begin{watchme}\begin{minipage}[c]{0.4\linewidth}\centering\threept{A}{B}{C}{$11$}{$?$}{$30$}\end{minipage}\hfill
\begin{minipage}[c]{0.56\linewidth}You know the whole and one piece, so subtract.\par
$BC = AC - AB = 30 - 11$. \quad \textbf{Answer: $\boldsymbol{BC = 19}$}\end{minipage}\end{watchme}
\needspace{3.4cm}\noindent\sdp{5}{$AC = 25$ and $BC = 9$. Find $AB$.}{$AB = 16$}\hfill\sdp{6}{$AC = 17.5$ and $AB = 6.8$. Find $BC$.}{$BC = 10.7$}\par\vspace{2mm}
\needspace{3.4cm}\noindent\sdp{7}{$Q$ is between $P$ and $R$. $PR = 52$ and $QR = 29$. Find $PQ$.}{$PQ = 23$}\hfill\figp{8}{Find $MN$.}{\threept{L}{M}{N}{$7$}{$?$}{$18$}}{$MN = 11$}\par

\needspace{7cm}\topicbar{Type 3 \quad Solve with Algebra}
\begin{watchme}$AB = 2x + 1,\ BC = 3x - 4,\ AC = 27.$ \ Find $x$, $AB$ and $BC$.\par\vspace{1mm}
$(2x + 1) + (3x - 4) = 27 \ \Rightarrow\ 5x - 3 = 27 \ \Rightarrow\ 5x = 30 \ \Rightarrow\ x = 6$.\par
Don't stop at $x$: \ $AB = 2(6) + 1 = 13$, \ $BC = 3(6) - 4 = 14$. \quad Check: $13 + 14 = 27$ \ding{51}\par
\textbf{Answer: $\boldsymbol{x = 6,\ AB = 13,\ BC = 14}$}\end{watchme}
\needspace{5cm}\noindent\algp{9}{$AB = x + 4,\ BC = 3x,\ AC = 24$. Find $x$ and $BC$.}{$x = 5,\ BC = 15$}\hfill\algp{10}{$K$ is between $J$ and $L$. $JK = 2x - 3,\ KL = x + 6,\ JL = 18$. Find $x$ and $JK$.}{$x = 5,\ JK = 7$}\par\vspace{2mm}
\needspace{5cm}\noindent\algp{11}{$E$ is between $D$ and $F$. $DE = 4x - 5,\ EF = 2x + 1,\ DF = 32$. Find $x$ and $EF$.}{$x = 6,\ EF = 13$}\hfill\algp{12}{$M$ is between $A$ and $B$. $AM = 3x + 2,\ MB = 5x - 6,\ AB = 36$. Find $x$ and $AM$.}{$x = 5,\ AM = 17$}\par

\needspace{6cm}\topicbar{Type 4 \quad Is It Between? \& Number Lines}
\begin{watchme}$AB = 4,\ BC = 6,\ AC = 9$. Is $B$ between $A$ and $C$?\par
If it were, then $AB + BC$ would equal $AC$. But $4 + 6 = 10 \ne 9$. \quad \textbf{Answer: no.}\par
\emph{Number line:} the distance between two points is the difference of their coordinates, like $|9 - 2| = 7$.\end{watchme}
\needspace{3.4cm}\noindent\sdp{13}{$AB = 7,\ BC = 5,\ AC = 12$. Is $B$ between $A$ and $C$? Why?}{yes: $7 + 5 = 12$}\hfill\sdp{14}{$AB = 8,\ BC = 3,\ AC = 10$. Is $B$ between $A$ and $C$? Why?}{no: $8 + 3 = 11 \ne 10$}\par\vspace{2mm}
\needspace{3.8cm}\noindent\begin{minipage}[t][3.2cm]{0.48\linewidth}{\color{fmturn}\bfseries 15.}\ On a number line, $A$ is at $-3$, $B$ is at $2$ and $C$ is at $9$. Find $AB$, $BC$ and $AC$.\par\vfill\hfill\pans[4.2cm]{$5,\ 7,\ 12$}\end{minipage}\hfill
\begin{minipage}[t][3.2cm]{0.48\linewidth}{\color{fmturn}\bfseries 16.} \LH\ $B$ is between $A$ and $C$, $AB : BC = 2 : 3$, and $AC = 30$. Find $AB$.\par\vfill\hfill\pans[3.6cm]{$AB = 12$}\end{minipage}\par
\end{document}
